% BTS CIEL 1 — Réseaux — 08 — DNS
% Compiler deux fois avec XeLaTeX depuis fichiers/.
\documentclass[aspectratio=169,11pt]{beamer}
\usepackage{fontspec}
\setmainfont{Latin Modern Sans}
\setsansfont{Latin Modern Sans}
\setmonofont{Latin Modern Mono}
\usepackage{polyglossia}
\setmainlanguage{french}
\usepackage{tikz}
\usetikzlibrary{arrows.meta,positioning,calc,shapes.geometric}
\usepackage{listings}
\usepackage{xcolor}
\definecolor{foret}{HTML}{1C322D}
\definecolor{creme}{HTML}{F8F3EE}
\definecolor{or}{HTML}{EBB552}
\definecolor{rose}{HTML}{F1CDBE}
\definecolor{vertclair}{HTML}{DDE9DF}
\definecolor{encre}{HTML}{24372E}
\setbeamercolor{normal text}{fg=creme,bg=foret}
\setbeamercolor{frametitle}{fg=or,bg=foret}
\setbeamercolor{title}{fg=creme}
\setbeamercolor{block title}{fg=foret,bg=or}
\setbeamercolor{block body}{fg=encre,bg=creme}
\setbeamercolor{alerted text}{fg=or}
\setbeamercolor{structure}{fg=or}
\setbeamertemplate{navigation symbols}{}
\setbeamertemplate{footline}{\hfill\color{or}\insertframenumber\hspace{0.5cm}\vspace{0.25cm}}
\setbeamertemplate{frametitle}{\vspace{0.35cm}\insertframetitle\par\vspace{0.15cm}\color{or}\rule{\linewidth}{0.6pt}}
\setlength{\parskip}{5pt}
\lstdefinestyle{webcode}{basicstyle=\ttfamily\small\color{encre},keywordstyle=\color{foret},stringstyle=\color{foret},
  backgroundcolor=\color{creme},frame=single,rulecolor=\color{or},framerule=1pt,
  xleftmargin=5pt,xrightmargin=5pt,columns=fullflexible,keepspaces=true,showstringspaces=false}
\newcommand{\question}[1]{\begin{beamercolorbox}[rounded=true,sep=12pt]{block body}\Large\textbf{#1}\end{beamercolorbox}}
\newcommand{\etiquette}[1]{{\color{or}\small\MakeUppercase{#1}}}
% Les schémas ont une surface claire, également lisible en noir et blanc.
\tikzset{
  boite/.style={draw=encre,fill=creme,text=encre,rounded corners,
    align=center,text width=2.8cm,minimum height=1cm},
  flux/.style={-{Stealth},thick,draw=encre},
  every node/.style={text=encre,font=\small}}
\newenvironment{schema}{%
  \begin{beamercolorbox}[sep=8pt,rounded=true]{block body}
  \centering\begin{tikzpicture}[>=Stealth]
}{\end{tikzpicture}\end{beamercolorbox}}
\title{Le protocole DNS}
\subtitle{BTS CIEL 1 · Réseaux · 08 — DNS}
\author{}
\date{}
\begin{document}

\begin{frame}[plain]
  \vfill
  \etiquette{Retrouver une adresse à partir d'un nom}\par\bigskip
  {\Huge\bfseries Le protocole DNS}\par\bigskip
  {\Large BTS CIEL 1 · Réseaux · 08 — DNS}\par\bigskip
  \begin{schema}
    \node[boite,text width=4.5cm,minimum height=1.5cm,font=\large] (n) {\texttt{www.example.com}};
    \node[boite,text width=4cm,minimum height=1.5cm,font=\large,right=3cm of n] (i) {\texttt{203.0.113.80}};
    \draw[flux] (n)--node[above]{DNS} (i);
  \end{schema}
  \vfill
\end{frame}

\begin{frame}{Un nom facile à retenir}
  \question{Comment appeler sans retenir un numéro ?}
  \begin{schema}
    \node[boite,text width=3cm,minimum height=1.5cm,font=\normalsize] (a) at (0,1.5) {Nom dans le répertoire\\« Maman »};
    \node[boite,text width=3cm,minimum height=1.5cm,font=\normalsize] (b) at (7,1.5) {Numéro de téléphone};
    \draw[flux] (a)--node[above,font=\small]{chercher} (b);
    \node[boite,text width=3.2cm,minimum height=1.5cm,font=\normalsize] (c) at (0,-.5) {Nom de domaine\\\texttt{www.example.com}};
    \node[boite,text width=3cm,minimum height=1.5cm,font=\normalsize] (d) at (6.9,-.5) {Adresse IP\\\texttt{203.0.113.80}};
    \draw[flux] (c)--node[above,font=\small]{rechercher} (d);
  \end{schema}
  {\small Les machines se joignent par IP. Sans annuaire, il faudrait retenir \texttt{203.0.113.80}.}
\end{frame}

\begin{frame}{Le navigateur demande d'abord l'adresse}
  Exemple fictif : \texttt{http://www.example.com}.
  \begin{schema}
    \node[font=\normalsize] at (0,3.5) {Navigateur};
    \node[font=\normalsize] at (5,3.5) {Serveur DNS};
    \node[font=\normalsize,align=center] at (10,3.5) {Serveur HTTP\\\texttt{203.0.113.80}};
    \draw[dashed] (0,3)--(0,.3); \draw[dashed] (5,3)--(5,1.3); \draw[dashed] (10,3)--(10,.3);
    \draw[flux] (0,2.7)--node[above,fill=creme,inner sep=2pt,font=\scriptsize]{1. IPv4 de www.example.com ?} (5,2.7);
    \draw[flux,dashed] (5,1.8)--node[above,fill=creme,inner sep=2pt,font=\small]{2. 203.0.113.80} (0,1.8);
    \draw[flux] (0,.6)--node[above,fill=creme,inner sep=2pt,font=\small]{3. Requête HTTP vers 203.0.113.80} (10,.6);
  \end{schema}
  Cette question-réponse s'appelle la résolution de nom.
\end{frame}

\begin{frame}{L'enregistrement A : une ligne de l'annuaire}
  Sur le serveur DNS, un enregistrement A associe un nom à une adresse IPv4.
  \begin{beamercolorbox}[sep=12pt,rounded=true]{block body}
    \centering\renewcommand{\arraystretch}{1.6}
    \begin{tabular}{lll}
      \textbf{Nom} & \textbf{Type} & \textbf{Valeur}\\\hline
      \texttt{www.example.com} & A & \texttt{203.0.113.80}\\
    \end{tabular}
  \end{beamercolorbox}
  \begin{block}{Méthode}
    1. Repérez le nom demandé.\\
    2. Trouvez la ligne de type A de ce nom.\\
    3. Lisez l'adresse IPv4.
  \end{block}
  Le type AAAA associe un nom à une adresse IPv6.
\end{frame}

\begin{frame}{Comment joindre le serveur DNS ?}
  Le poste doit connaître l'adresse du serveur DNS.
  \begin{schema}
    \node[boite,text width=2.8cm,minimum height=2.5cm,font=\large] (d) at (0,0) {DHCP ou\\saisie manuelle};
    \node[boite,text width=2.7cm,minimum height=2.5cm,font=\large] (p) at (5,0) {Poste client};
    \node[boite,text width=2.8cm,minimum height=2.5cm,font=\large] (s) at (10,0) {Serveur DNS\\\texttt{192.168.10.53}};
    \draw[flux] (d)--node[above=1mm,font=\small,fill=creme,inner sep=2pt]{1. adresse} (p);
    \draw[flux] (p)--node[above=1mm,font=\small,fill=creme,inner sep=2pt]{2. question} (s);
  \end{schema}
  Les requêtes DNS utilisent généralement UDP, port serveur \textbf{53}.
\end{frame}

\begin{frame}{Les noms forment un arbre}
  Lisez le nom de droite à gauche pour parcourir la hiérarchie.
  \begin{schema}
    \node[boite,text width=2cm,font=\normalsize] (r) at (5,4) {Racine : \texttt{.}};
    \node[boite,text width=1.5cm,font=\normalsize] (com) at (2,2.8) {\texttt{com}};
    \node[boite,text width=1.5cm,font=\normalsize] (org) at (5,2.8) {\texttt{org}};
    \node[boite,text width=1.5cm,font=\normalsize] (fr) at (8,2.8) {\texttt{fr}};
    \node[boite,text width=3cm,font=\normalsize] (ex) at (2,1.5) {\texttt{example.com}};
    \node[boite,text width=3.6cm,font=\normalsize] (www) at (0,0) {\texttt{www.example.com}};
    \node[boite,text width=3.6cm,font=\normalsize] (mail) at (4.3,0) {\texttt{mail.example.com}};
    \draw (r)--(com); \draw (r)--(org); \draw (r)--(fr); \draw (com)--(ex); \draw (ex)--(www); \draw (ex)--(mail);
  \end{schema}
\end{frame}

\begin{frame}{Trouver une réponse absente du cache}
  Le cache du serveur DNS local conserve temporairement les réponses. Ici, la réponse est absente.
  \begin{schema}
    \node[font=\normalsize] at (0,4) {Poste};
    \node[font=\normalsize,align=center] at (4,4.2) {Serveur DNS local\\\texttt{192.168.10.53}};
    \node[font=\normalsize] at (9,4) {Serveur racine};
    \foreach \x in {0,4,9} \draw[dashed] (\x,3.5)--(\x,.5);
    \draw[flux] (0,3)--node[above,fill=creme,inner sep=2pt,font=\scriptsize]{1. IPv4 de www.example.com ?} (4,3);
    \draw[flux] (4,2)--node[above,fill=creme,inner sep=2pt,font=\scriptsize]{2. IPv4 de www.example.com ?} (9,2);
    \draw[flux,dashed] (9,.8)--node[above,fill=creme,inner sep=2pt,font=\small]{3. adresse d’un serveur .com} (4,.8);
  \end{schema}
\end{frame}
\begin{frame}{Poursuivre la recherche}
  \begin{schema}
    \node[font=\normalsize,align=center] at (0,4.5) {Serveur DNS local\\\texttt{192.168.10.53}};
    \node[font=\normalsize,align=center] at (5,4.5) {Serveur\\\texttt{.com}};
    \node[font=\normalsize,align=center] at (10,4.5) {Serveur\\\texttt{example.com}};
    \draw[dashed] (0,3.9)--(0,.8);
    \draw[dashed] (5,3.9)--(5,2.5);
    \draw[dashed] (10,3.9)--(10,.8);
    \draw[flux] (0,3.5)--node[above,fill=creme,inner sep=2pt,font=\scriptsize]{4. IPv4 de www.example.com ?} (5,3.5);
    \draw[flux,dashed] (5,2.7)--node[above,fill=creme,inner sep=2pt,font=\scriptsize]{5. adresse du serveur de example.com} (0,2.7);
    \draw[flux] (0,1.9)--node[above,fill=creme,inner sep=2pt,font=\small]{6. IPv4 de www.example.com ?} (10,1.9);
    \draw[flux,dashed] (10,1.1)--node[above,fill=creme,inner sep=2pt,font=\small]{7. réponse A} (0,1.1);
  \end{schema}
  8. Il garde la réponse en cache et la transmet au poste.
\end{frame}

\begin{frame}{Diagnostiquer un site qui ne s'ouvre pas par son nom}
  \begin{block}{Méthode}
    1. Testez l'IP : Packet Tracer ou Windows, \texttt{ping 203.0.113.80}.\\
    2. Vérifiez le DNS du poste : Packet Tracer, PC > Desktop > IP Configuration, champ DNS Server.\\
    3. Vérifiez l'enregistrement A : Packet Tracer, serveur > Services > DNS.
  \end{block}
  \begin{schema}
    \node[boite,text width=3.6cm,minimum height=2cm,font=\normalsize] (a) at (0,0) {1. Test IP en échec :\\vérifiez le réseau};
    \node[boite,text width=3.6cm,minimum height=2cm,font=\normalsize] (b) at (4.7,0) {2. DNS incorrect :\\corrigez l'adresse};
    \node[boite,text width=3.6cm,minimum height=2cm,font=\normalsize] (c) at (9.4,0) {3. Ligne A incorrecte :\\corrigez-la};
    \draw[flux] (a)--(b); \draw[flux] (b)--(c);
  \end{schema}
\end{frame}

\begin{frame}{À vous : lire l'annuaire}
  \question{Quelle IPv4 pour \texttt{http://atelier.example.org} ? Justifiez.}
  \begin{beamercolorbox}[sep=12pt,rounded=true]{block body}
    \centering\renewcommand{\arraystretch}{1.5}
    \begin{tabular}{lll}
      \textbf{Nom} & \textbf{Type} & \textbf{Valeur}\\\hline
      \texttt{atelier.example.org} & AAAA & \texttt{2001:db8::20}\\
      \texttt{atelier.example.com} & A & \texttt{198.51.100.30}\\
      \texttt{atelier.example.org} & A & \texttt{198.51.100.20}\\
    \end{tabular}
  \end{beamercolorbox}
\end{frame}

\begin{frame}{Correction : lire la valeur A}
  \begin{beamercolorbox}[sep=12pt,rounded=true]{block body}
    \centering\renewcommand{\arraystretch}{1.8}
    \begin{tabular}{ll}
      \textbf{Critère} & \textbf{Choix}\\\hline
      Nom demandé & \texttt{atelier.example.org}\\
      IPv4 & type A, pas AAAA\\
      Adresse retenue & \texttt{198.51.100.20}\\
    \end{tabular}
  \end{beamercolorbox}
  Le nom en \texttt{.com} est différent. AAAA désigne une IPv6.
\end{frame}

\begin{frame}{À vous : trouver le serveur DNS}
  \question{Quelle étape de diagnostic trouve la panne ? Pourquoi ?}
  \begin{schema}
    \node[boite,text width=3.2cm,minimum height=1.7cm] (p) at (0,1) {Poste\\DNS configuré :\\\texttt{192.168.10.54}};
    \node[boite,text width=3.2cm,minimum height=1.7cm] (x) at (5,1) {\texttt{192.168.10.54}\\aucun serveur DNS};
    \node[boite,text width=3.2cm,minimum height=1.7cm] (d) at (10,1) {Serveur DNS\\\texttt{192.168.10.53}};
    \draw[flux,dashed] (p)--node[above,font=\small]{question} (x);
    \node[font=\normalsize] at (5,-.7) {\texttt{ping 203.0.113.80} (serveur Web) répond.};
  \end{schema}
\end{frame}

\begin{frame}{Correction : la bonne adresse DNS}
  \begin{schema}
    \node[boite,text width=3.2cm,minimum height=1.7cm] (p) at (0,1) {Poste\\DNS configuré :\\\texttt{192.168.10.54}};
    \node[boite,text width=3.2cm,minimum height=1.7cm] (x) at (5,1) {\texttt{192.168.10.54}\\aucun serveur DNS};
    \node[boite,text width=3.2cm,minimum height=1.7cm] (d) at (10,1) {Serveur DNS\\\texttt{192.168.10.53}};
    \draw[flux,dashed] (p)--node[above,font=\small]{question} (x);
    \node[font=\normalsize] at (5,-.7) {\texttt{192.168.10.54} $\neq$ \texttt{192.168.10.53}};
  \end{schema}
  \begin{block}{Étape 2 : vérifier le DNS du poste}
    Le test IP répond. L'adresse DNS configurée ne correspond pas au serveur.\\
    Corrigez l'adresse DNS du poste en \texttt{192.168.10.53}.
  \end{block}
\end{frame}

\begin{frame}{À retenir}
  \begin{schema}
    \node[boite,text width=3.8cm,minimum height=1.8cm,font=\large] (n) {Nom de domaine};
    \node[boite,text width=3.8cm,minimum height=1.8cm,font=\large,right=1cm of n] (d) {Serveur DNS\\enregistrement A};
    \node[boite,text width=3cm,minimum height=1.8cm,font=\large,right=1cm of d] (i) {Adresse IPv4};
    \draw[flux] (n)--(d); \draw[flux] (d)--(i);
  \end{schema}
  \begin{block}{Quatre repères}
    DNS réalise la résolution de nom.\\
    Un enregistrement A associe un nom à une IPv4.\\
    L'adresse DNS du poste est saisie ou reçue par DHCP.\\
    Les requêtes utilisent généralement UDP, port serveur 53.
  \end{block}
\end{frame}

\begin{frame}{Sources}
  \begin{beamercolorbox}[sep=12pt,rounded=true]{block body}
    \small
    \textbf{RFC 1034} — concepts et hiérarchie DNS\\
    \texttt{https://www.rfc-editor.org/rfc/rfc1034}\\[4pt]
    \textbf{RFC 1035} — messages, enregistrements et transport\\
    \texttt{https://www.rfc-editor.org/rfc/rfc1035}\\[4pt]
    \textbf{RFC 3596} — enregistrements AAAA pour IPv6\\
    \texttt{https://www.rfc-editor.org/rfc/rfc3596}\\[4pt]
    \textbf{RFC 2132} — adresse DNS fournie par DHCP, section 3.8\\
    \texttt{https://www.rfc-editor.org/rfc/rfc2132}\\[4pt]
    Support enseignant : « TP DHCP, HTTP, DNS, FTP », p. 1--3 et 9--11.
  \end{beamercolorbox}
\end{frame}
\end{document}
